\section{Phonon-assisted electron recoil from first principles}\label{app:E.14}

The factor-of-ten estimate of the phonon-assisted rate of silicon (\cref{sec:13.9}) takes the measured indirect absorption of silicon for the optical limit and multiplies it by a computed form factor of the zero-phonon vertex for the fall-off at finite $q$. It never evaluates the four-path amplitude of \cref{sec:13.2}. The amplitude is hard for three reasons: the band sum over intermediate states is unrestricted and cannot stop at the top of a Wannier manifold; the density vertex is needed at momentum transfers of one to a few keV, where the small-$q$ expansion in Wannier position matrix elements fails (\cref{app:E.5}); and the external screening carries local fields, which the macroscopic dielectric function does not. Everything else follows the standard steps of \cref{app:E.5}.

\subsection{The phonon-assisted loss-function grid}

What a rate needs is the phonon-assisted loss function $L^{\mathrm{PA}}(\vect{q},\omega)$ of \cref{sec:13.3}, a dimensionless number per bin of $(|\vect{q}|,\hat{\vect{q}},\omega)$, built from the four-path amplitude of \cref{sec:13.2} and screened as in \cref{sec:13.4}. The bins are those of the estimate in \cref{sec:13.9}, so that the two can be compared bin by bin. They are $|\vect{q}| = 0.25$ to \SI{1.25}{\angstrom^{-1}} (0.49 to \SI{2.47}{keV}) in steps of \SI{0.05}{\angstrom^{-1}}, 30 directions in the irreducible wedge of the cubic group, and $\omega$ from $\Eg + \omega_{\mathrm{TA}}(\vect{Q}^*)$ up to the pointwise window edge $E_{\mathrm{min}}(\vect{q}) - \delta_w$ of \cref{sec:13.6}. From it follow the rate $\dd\mathcal{R}/\dd\omega$ and the reach of \cref{sec:13.3}, for the heavy and the light mediator. On the truncation label $\mathfrak{T}$ (\cref{sec:8.5}) the calculation sits at $n_B = 0$, $n_g = 2$ (one density vertex and one electron--phonon vertex, all four time orderings), $n_T = 0$, lattice regime $\mathcal{K}$ = one phonon, and $\mathcal{X}$ = RPA with local fields for the external screening. The free-pair amplitude has no weight below $\Eg + \omega_{\mathrm{TA}}$; the replica shell below it (0.3332 to \SI{0.3374}{MeV} in mass) is the exciton envelope times the same amplitude, matched at $\Eg + \omega_{\mathrm{TA}}$.

The ground state and the phonons are the shared harmonic base of \cref{app:E.17}, with its $12^3$ $\vect{k}$ grid and $8^3$ wavevector grid. On a $4^3$ grid the interpolated transverse-acoustic frequency at $0.85\,\mathrm{X}$ comes out above its value at X, an artifact of up to \SI{1.5}{meV} on exactly the branch that sets the window. The self-consistent potential changes are kept for every wavevector, since the vertices are built from them, and the long-range quadrupolar part of $g$ is split off before interpolation. An 18-function Wannier manifold ($sp^3$ on each atom plus $d$-like functions) is disentangled from about 60 bands on the full $12^3$ grid without symmetry. The conduction states to \SI{10}{eV} above the gap then sit inside the frozen window, and the size of the manifold is tested by repeating the amplitude with 26 functions.

Since the band sum multiplies matrices in band space, $g$, the density vertex $\rho_{mn}(\vect{k},\vect{q})$ and the energies have to be rotated with the same $U(\vect{k})$ (\cref{app:A.5}); with a mismatch every band energy is still correct but the phases of the sum are wrong. The acoustic sum rule is imposed on the dynamical matrix, never by setting $g$ to zero as $\vect{Q}\to0$, because the part of $g$ linear in $\vect{Q}$ is physical and is exactly what the soft acoustic limit of \cref{sec:13.6} tests.

The density vertex is computed for each target $\vect{q}$ from the plane-wave coefficients of one non-self-consistent run at the shifted points $\vect{k}+\vect{q}$, with $\vect{q}$ kept unreduced (\cref{app:A.3}); a vertex at a reduced $\vect{q}$ belongs to a different Umklapp class. With 21 values of $|\vect{q}|$ and 30 directions there are 630 targets in the irreducible wedge, about 1800 once equivalent pockets are unfolded, and each needs the full $12^3$ grid without symmetry, about $10^6$ diagonalizations in all. The amplitude is then the four-path sum with an unrestricted band sum: explicit inside the Wannier manifold, and above it the projected Sternheimer solution of \cref{eq:E.4-dpsiperp} contracted with $\ee^{\ii\vect{q}\cdot\vect{r}}$, one linear solve per external state and perturbation. That remainder matters most on the Bethe ridge, where $q^2/2m_e$ approaches the energy deposited. Below $|\vect{q}| \approx 0.5\,\alpha m_e \approx \SI{1.9}{keV}$ it is a correction to the Wannier sum; on the upper bins, 1.9 to \SI{2.5}{keV}, it is carried in full. The external screening is the column $(\tilde\varepsilon^{-1})_{\vect{G}0}(\vect{q},\omega)$ of an RPA calculation with local fields (\cref{app:E.6}). The pair channel of \cref{sec:13.4} adds $W^{\mathrm{scr}}(\vect{q}+\vect{Q},\omega_f)$ at the pair momentum $\vect{q}+\vect{Q}$, which kinematics pins near $0.85\,\mathrm{X}$, so one dielectric matrix serves every bin. The two corrections are reported as the ratios $\delta_{\mathrm{LF}}$ and $\delta_{\mathrm{pair}}$ of \cref{sec:13.4}. The sums over pocket pairs run on fine grids of $60^3$, refined to $120^3$, with tetrahedra and no broadening.

Quantum ESPRESSO~\cite{giannozzi2017advanced}, Wannier90~\cite{pizzi2020wannier90}, EPW~\cite{ponce2016epw} and BerkeleyGW~\cite{deslippe2012berkeleygw} carry most of this, and a long-wave DFPT code the quadrupoles. The density vertex in the Wannier gauge and the second-order amplitude with its Sternheimer remainder are in none of them; they are new code, a few hundred lines on top of EPW. \Cref{tab:E14-params} collects parameters and costs. The plane-wave stages total about 350 to 450 core-hours, dominated by the density vertices; the amplitude, at the $10^{16}$ operations that \cref{app:E.5} counts, adds a quarter of an hour to a few hours on one accelerator.

\begin{table}[t]
  \centering
  \caption{Parameters and costs of the phonon-assisted loss function of silicon. The cutoff refers to the SG15 norm-conserving pseudopotential~\cite{schlipf2015optimization} with four valence electrons. Fine grids are those of the pockets within \SI{0.1}{eV} of the band extrema, quoted as their equivalent uniform density. Costs are in core-hours unless stated; ``base'' is the shared DFPT run of \cref{app:E.17}.}
  \label{tab:E14-params}
  \small
  \setlength{\tabcolsep}{4pt}
  \begin{tabular}{p{0.16\textwidth}p{0.28\textwidth}p{0.3\textwidth}p{0.14\textwidth}}
    \toprule
    stage & grids and parameters & acceptance test & cost \\
    \midrule
    ground state, phonons & PBE, $a = \SI{5.431}{\angstrom}$, \SI{60}{Ry}, $12^3$ $\vect{k}$, $8^3$ $\vect{Q}$ (29 irreducible) & neutron data per branch (\cref{app:E.17}) & base, $\approx 100$ \\
    long-range $g$ & dipoles (zero in Si), quadrupoles & small-$\vect{Q}$ angular dependence & minutes \\
    Wannier gauge & 18 functions ($sp^3$, $d$); window to \SI{40}{eV} & bands $\le\SI{10}{meV}$ to \SI{10}{eV} above the gap & 3--5 \\
    density vertex & 630 targets, $12^3$ without symmetry & direct vs Wannier at 5 off-grid $\vect{k}$, $\le\SI{2}{\percent}$ & 200--300 \\
    vertex $g$ & interpolated & direct DFPT at 3 off-grid $\vect{Q}$, $\le\SI{5}{\percent}$ & 10 \\
    amplitude & $60^3 \to 120^3$, tetrahedra; Wannier sum plus Sternheimer remainder ($10^3$ solves per target) & \SI{5}{\percent} per bin in $L^{\mathrm{PA}}$ & accelerator hours (\cref{app:E.5}) \\
    screening & RPA, $(\tilde\varepsilon^{-1})_{\vect{G}0}(\vect{q},\omega)$; $W^{\mathrm{scr}}$ at $\vect{q}+\vect{Q}$ & $\delta_{\mathrm{LF}}$, $\delta_{\mathrm{pair}}$ reported & tens \\
    \bottomrule
  \end{tabular}
\end{table}

\subsection{Physical controls}

A bin converged to \SI{5}{\percent} under the refinements of \cref{tab:E14-params} need not be correct, and each control below tests the physics against something the calculation was not fitted to.

The optical limit tests the vertices against data. At $|\vect{q}| = \SI{e-3}{\angstrom^{-1}}$ the density vertex reduces to its dipole form (\cref{sec:3.2}), and $L^{\mathrm{PA}}$ has to reproduce the phonon-assisted $\varepsilon_2^{\mathrm{PA}}$ that EPW computes from the same vertices. After the gap shift it also has to reproduce the measured absorption edge at the temperature of each measurement: silicon at \SI{20}{K}, where the transverse-acoustic and transverse-optical replicas are resolved~\cite{macfarlane1958fine}, germanium against its edge data at 77 to \SI{80}{K}, and GaAs at the temperature of its data. The comparison is made branch by branch, with $\varepsilon_2^{\mathrm{PA}}$ resolved by phonon branch, since the transfer of a measured edge to millikelvin needs the branch weights and a measured spectrum gives only their sum (\cref{sec:13.5}). Switching $g$ off tests the labeling of final states: the phonon-assisted rate then vanishes identically, and any residue is zero-phonon weight leaking into the window through a mislabeled state.

A frozen phonon tests the adiabatic form of the amplitude derived in \cref{sec:13.2}. At $\vect{Q} = \mathrm{X}$ an eight-atom cell holds the phonon exactly; displacing the atoms by $\pm\SI{0.005}{\angstrom}$ along a mode and differencing the zero-phonon transition densities gives the adiabatic amplitude directly, which agrees with the second-order sum up to the expected $\omega_{\vect{Q}\nu}/\delta_{\mathrm{off}}$. The soft acoustic limit of \cref{sec:13.6}, checked at $|\vect{Q}| = 10^{-2}$, $10^{-3}$ and $10^{-4}$ of $2\pi/a$ in three directions, fails at once for a vertex whose small-$\vect{Q}$ part was zeroed. Symmetry gives two exact zeros: on the valley axis the longitudinal-acoustic path of the class without Umklapp vanishes, and so does the longitudinal-optical path of the $(\bar2\,0\,0)$ class (\cref{sec:13.9}); both come out at the level of numerical noise.

Between 0.333 and about \SI{0.35}{MeV} the rate is a power of the distance from the onset (\cref{sec:13.3}), so there the exponent is the meaningful output.

\subsection{Zero-phonon rates on common bins}

\Cref{sec:13.9} places the end of the phonon-assisted lead at about 0.42 to \SI{0.43}{MeV}, with a two-channel band to about \SI{0.51}{MeV}. That is a ratio of two rates, settled only when both are computed on the same bins, so the output grid extends past the window into three regions. The first zero-phonon edge, 0.405 to \SI{0.408}{MeV} with a best value of \SI{0.407}{MeV}, is a joint minimum over the momentum transfer and the hole's crystal momentum and lies off the valley axis. Its state belongs to the symmetry-suppressed (glide) class, so its vertex is small and it does not end the lead. Around this mirror-plane optimum the bins are cones: $|\vect{q}| = 1.2$ to \SI{1.9}{keV}, tilts $\vartheta \le \SI{0.4}{rad}$ from the axis within the $\{001\}$ and $\{011\}$ planes and up to \SI{0.15}{rad} out of plane, and $\omega$ within \SI{0.3}{eV} of $E_{\mathrm{min}}(\vect{q})$. The second region is the cone of the second heavy hole on the axis, the mirror-even band that opens from 0.415 to \SI{0.420}{MeV}. The third is the X-face edge of the allowed $(\bar2\,0\,0)$ class, where the zero-phonon line of allowed size opens near \SI{0.50}{MeV}. Above $E_{\mathrm{min}}(\vect{q})$ the phonon-assisted value follows the convention of \cref{sec:13.6}; the shell just below the edge is reported separately, and the bins contain the joint optimum at \SI{0.407}{MeV}, so that the edge is the one the rate sees.

The zero-phonon side costs under one core-hour. It is a Wannier--Mott continuum (\cref{sec:4.9}) with the anisotropic reduced-mass tensor of the pair at each bin, from the Hessians of the interpolated bands at the pair's two crystal momenta. It uses a static $\varepsilon_\infty$ and density vertices computed directly from the plane-wave states of both top heavy-hole bands. The hole is placed where the pair is lowest at that $\vect{q}$; at the off-axis optimum it sits about \SI{0.06}{eV} below the valence maximum, and placing it at $\mathrm{\Gamma}$ misses the optimum. At the class extremum the allowed $(\bar2\,0\,0)$ vertex is $\sum_v|\rho_{cv}|^2 = 0.150$, the size of an ordinary allowed vertex (\cref{sec:13.9}); the comparison needs it on every bin, with spin--orbit coupling in the bands, which shifts the first-hole bins by about \SI{15}{meV}. The vertices come from matching Miller indices on a three-point stencil around each bin. They are converged when halving the stencil step from $0.02$ to $0.01\times2\pi/a$ changes $|\rho|^2$ by less than \SI{2}{\percent}. Halving the angular step from 0.01 to \SI{0.005}{rad} has to change the rate by less than \SI{1}{\percent}, and the ratio of the two rates per bin has to be stable to \SI{5}{\percent} when the number of bins grows by half. Off the axis the band index no longer orders the states by mirror parity, and each state is assigned by its parity; on the axis the degenerate pair is summed.

\subsection{Germanium and gallium arsenide}

Germanium cannot follow silicon, because PBE gets its band order wrong. It puts the conduction state at $\mathrm{\Gamma}$ \SI{0.047}{eV} above the valence maximum and the L valley above it, the reverse of the measured order, in which the indirect gap at L (\SI{0.744}{eV}) lies below the direct gap at $\mathrm{\Gamma}$ (\SI{0.89}{eV}). A scissor applied to that order gives a window that is wrong by the inversion of $\mathrm{\Gamma}$ and L. Any germanium band structure is therefore checked before it is used: L below $\mathrm{\Gamma}$, both gaps within \SI{0.1}{eV} of the measured ones, and the band masses at both valleys against measurement, the transverse mass at L within $0.010\,m_e$ of $m^*_\perp = 0.0815\,m_e$ (\cref{tab:3-materials}) and the mass at $\mathrm{\Gamma}$ within $0.005\,m_e$ of $0.038\,m_e$.

HSE06 with spin--orbit coupling at the standard exchange fraction of one quarter fails this check. At the measured lattice constant it puts $\mathrm{\Gamma}$ at 0.64 to \SI{0.69}{eV} and L at 0.79 to \SI{0.83}{eV} above the valence maximum, so the gap at L passes while the order and the direct gap fail, with an inversion of 0.09 to \SI{0.15}{eV}. Above a fraction of 0.28 the order is restored, but at 0.3 the direct and indirect gaps are still \SI{34}{\percent} and \SI{17}{\percent} off~\cite[Fig.~1(a)]{abdurrazaq2025accurate}. A correction that shifts only the 4s-like orbitals~\cite{abdurrazaq2025accurate}, with its single parameter fitted to both gaps, passes the energies, L at \SI{0.749}{eV} below $\mathrm{\Gamma}$ at \SI{0.918}{eV}. It fails the masses: the transverse mass at L comes out $0.101\,m_e$ and the mass at $\mathrm{\Gamma}$ $0.053\,m_e$, \SI{24}{\percent} and \SI{39}{\percent} heavy, so the valley energies are right and the curvatures are not. The next route is a quasiparticle $GW$ correction (\cref{app:E.3}) on top of the hybrid or PBE bands, subject to the same check. Until bands pass, the zero-phonon edge of germanium stays the estimate of \cref{sec:13.1}, \SI{0.32}{MeV} on the L axis inside a bracket of 0.212 to \SI{0.420}{MeV}.

On bands that pass, sixteen spinor $sp^3$ Wannier functions are extracted with the twenty 3d bands excluded and a frozen window to \SI{3}{eV} above the valence maximum, and the interpolated edges and masses are compared with direct calculations at off-grid points. The zero-phonon edge then follows by the joint minimization over $(\vect{q},\vect{k})$ for the L-valley class and its Umklapp copies. Whether that optimum leaves the axis is not known: the germanium heavy hole is less warped than that of silicon, and nothing transfers from the silicon result. The mirror rule of \cref{sec:13.9} becomes one for the diagonal mirror planes, which carry no glide phase; whether the first state is allowed is checked at the true optimum.

The germanium phonons are the shared base, with an $8^3$ $\vect{k}$ grid and a $6^3$ wavevector grid. On a $4^3$ wavevector grid the transverse-acoustic frequencies are \SI{13.5}{\percent} (X) and \SI{11.6}{\percent} (L) below neutron data at \SI{80}{K}~\cite{nilsson1971phonon}, while the optical and longitudinal branches agree to \SI{2.3}{\percent}. Either the nearly closed PBE gap overscreens the shear modes or the $8^3$ $\vect{k}$ grid is too coarse for a \SI{47}{meV} gap. One DFPT run at X on a $12^3$ grid, about three core-hours, separates the two, and $\omega_{\mathrm{TA}}(\vect{Q}^*)$ has to be stable to \SI{0.3}{meV} between the $4^3$ and $6^3$ grids. If the deficit survives, the neutron-scaled branch is used for the window; the phonon-assisted edge moves by at most \SI{0.0003}{MeV} either way. The germanium rate then follows the silicon chain with spin--orbit coupling and the 3d shell in the valence (14 electrons per atom): 36 spinor Wannier functions, the sixteen $sp^3$ ones plus twenty for the 3d shell, and 120 spinor bands. Spinors double both the length of each state and the band count, which multiplies each plane-wave stage by four to eight; for germanium's smaller DFPT grid and targets the chain comes to 400 to 800 core-hours, or hours on one accelerator node.

Gallium arsenide is the direct-gap control. There the zero-phonon pair at small $q$ is out of reach for any mass: the pair at $\mathrm{\Gamma}$ needs $E_{\mathrm{dir}} + q^2/2M_x \le qv_{\mathrm{max}}$, and the right side minus the recoil is at most $M_xv_{\mathrm{max}}^2/2 \approx \SI{1.04}{eV}$ for a pair mass $M_x \approx 0.58\,m_e$, below the direct gap of \SI{1.52}{eV}. A phonon can absorb the momentum instead, and the phonon-assisted rate leads until the L-valley zero-phonon state opens, from about \SI{0.43}{MeV} to between 0.52 and \SI{0.56}{MeV} in a parabolic estimate with the L valley at \SI{1.815}{eV} (\cref{sec:13.1,sec:13.9}). The chain, with the Fr\"ohlich vertex (\cref{eq:E.5-gL}), has to reproduce that band and give a small correction above the zero-phonon threshold; without the long-range part the vertex diverges at small $\vect{Q}$. The ground state is the shared base at $a = \SI{5.653}{\angstrom}$ with the gallium 3d shell in the valence. PBE puts the direct gap at about \SI{0.5}{eV}, so the bands come from HSE06 or from a $\vect{k}$-dependent scissor fitted to the low-temperature conduction minima at $\mathrm{\Gamma}$ (\SI{1.519}{eV}), L (\SI{1.815}{eV}) and X (\SI{1.981}{eV})~\cite{vurgaftman2001band}, with spin--orbit coupling for the holes. The Fr\"ohlich vertex takes $\varepsilon_\infty$ and $Z^*$ from measurement through \cref{eq:E.17-lst}, since the PBE $\varepsilon_\infty$ of GaAs is tens of percent too large. The window costs about 8 core-hours with the scissor and 40 to 80 with the hybrid; the full rate costs about as much as silicon's.

For the other indirect-gap candidates the window comes before any rate: $E_{\mathrm{min}}(\vect{q})$ on the joint $(\vect{q},\vect{k})$ space, the band masses, $\omega_{\mathrm{TA}}(\vect{Q}^*)$ and $g$ at $\vect{Q}^*$ from one DFPT run. PBE with a scissor serves wherever PBE gets the band order right, checked against measured gaps at the high-symmetry points; narrow-gap and d-electron crystals need bands that pass the germanium check of order, gaps and masses. The kinematics are checked in every case by the limit $|\vect{q}| \to |\vect{K}_{\mathrm{pk}}|$, the crystal momentum of the conduction pocket, where $E_{\mathrm{min}}$ returns the gap.

\subsection{Fano widths and detunings in GaAs}

Near the edge of the off-shell window of \cref{sec:13.6}, a zero-phonon final state of energy $\tilde E$ lies inside the phonon-assisted continuum and interferes with it, the Fano term of \cref{sec:8.6}. The term vanishes in crystals invariant under inversion combined with time reversal with an even electron number, such as silicon and germanium, at the level of one-body self-energies, but not in GaAs. Level shifts from the electron--hole vertex can split a degenerate zero-phonon manifold in a basis that is not invariant, and then only the sum of the Fano terms over the branches vanishes. The Fano term is perturbative only where the width $\varGamma$ of the resonance is small compared with its detuning $\delta_{\mathrm{F}}$, the distance of $\tilde E$ to the nearest feature of the continuum. With the Fano coefficient $c^F_\nu$ of each branch $\nu$, the branch probability of the first emission is $(\varGamma^{(\nu)}/\varGamma^{\mathrm{tot}})(1 + \bar c^F - c^F_\nu)$, with $\bar c^F$ the width-weighted mean. In GaAs $|c^F_\nu| \lesssim 0.3\varGamma/\delta_{\mathrm{F}}$ to first order, which bounds the change of the pair-number distribution by $0.3\varGamma/\delta_{\mathrm{F}}$ and keeps the probabilities positive for $\varGamma/\delta_{\mathrm{F}} \le 0.5$. What \cref{sec:13.6} needs is therefore a map of $\varGamma/\delta_{\mathrm{F}}$ over $(q,\omega)$.

The width of the pair state is the sum of the electron and hole widths $\varGamma_n(\vect{k})$. These are on-shell Fan--Migdal widths,
\begin{equation*}
  \varGamma_n(\vect{k}) = -2\Imag\Sigma^{\mathrm{FM}}_{n\vect{k}}(\varepsilon_{n\vect{k}}) ,
\end{equation*}
computed at \SI{10}{K} (\cref{sec:6.3}), at $n_g = 2$ with a statically screened vertex. The detuning is the smallest of the distances to the renormalized edge $E_{\min}^{\mathrm{ren}}(q)$, to the nearest intervalley threshold of the electron (L, X) or the split-off threshold of the hole, and the local scale of the joint density of states. The coefficient $c^F_\nu$ is the non-resonant remainder of the same second-order amplitude, evaluated on the final shell. Estimates give a Fr\"ohlich LO width of a few meV from just above the conduction minimum to the $\mathrm{\Gamma}$--L onset (about \SI{3}{meV} at the onset, \SI{0.3}{eV} above the minimum, with $\alpha_F = 0.068$ and $\omega_{\mathrm{LO}} = \SI{36}{meV}$), and 20 to \SI{60}{meV} above the onset, so $\varGamma/\delta_{\mathrm{F}}$ is 0.01 to 0.2 away from thresholds and diverges at each onset. The map excludes shells of $\pm2\varGamma$ around the $\mathrm{\Gamma}$--L onset and the split-off onset (\SI{0.34}{eV} below the valence maximum), and the shell within $3\varGamma$ of the window edge; inside them only totals from the unexpanded Dyson form are quoted.

Because the split-off band sets the hole widths, the electrons need fully relativistic pseudopotentials, at \SI{90}{Ry} because the spinor calculation with the Ga 3d shell converges more slowly than the scalar base. Phonons, Born charges and $\varepsilon_\infty$ come from the shared base on a $6^3$ grid. The run uses an $8^3$ $\vect{k}$ grid, 16 spinor $sp^3$ Wannier functions with the 3d bands excluded, the polar vertex, and self-energies on a $100^3$ fine $\vect{Q}$ grid for the 13\,824 $\vect{k}$ points of the final shell. The widths converge to \SI{5}{\percent} from $80^3$ to $100^3$ and are stable as the smearing goes from 10 to \SI{2.5}{meV}; $c^F_\nu$ converges to \SI{10}{\percent} when 20 bands are added to the intermediate sum. Silicon, with eight $sp^3$ functions and a $60^3$ grid, has to return $c^F_\nu = 0$ to numerical precision, which tests the chain against the symmetry theorem of \cref{sec:8.6}. The self-energies cost about $10^4$ core-hours (30 to 60 accelerator-hours), silicon a tenth of this. The map fixes the $(q,\omega)$ region in which \cref{sec:13.6} quotes the Fano term for GaAs.

\subsection{The phonon-assisted kinetic-energy row}

A scalar coupled to the electron's mass is absorbed below the direct gap of silicon only with a phonon, and its rate is the structure factor $S^{\mathrm{PA}}_{\mathcal{T}\mathcal{T}}(0,m_\phi)$ of \cref{eq:13.8-scalar}, for a point group without a polar vector (\cref{sec:10.7-scalar,sec:13.8}). No measurement supplies it, since its amplitude \cref{eq:13.8-MT} contains no velocity vertex and is not proportional to $\varepsilon_2^{\mathrm{PA}}$. The chain above computes it with one change: the density vertex at $\vect{q}$ is replaced by the kinetic vertex $\langle m|p^2/2m_e|n\rangle$ at $\vect{q} = 0$, and the energy is $\omega = m_\phi$ from the indirect edge to the direct gap. The Wannier gauge, the vertices $g$, the Sternheimer remainder above the manifold and the fine pocket grids are those of the density row. The kinetic matrix element converges like the total kinetic energy and is weighted toward the cores, so it is taken from PAW or all-electron states, at a cutoff of at least \SI{80}{Ry} for the norm-conserving comparison. Since $\langle c|p^2/2m_e|v\rangle = -\langle c|V|v\rangle$ between eigenstates of different energy, \cref{eq:13.8-MT} can be evaluated twice, once with the kinetic vertex and once with the matrix elements of the crystal potential, and the two must agree to the \SI{5}{\percent} per bin of \cref{tab:E14-params}; a disagreement points to the core region or to an incomplete band sum. The same chain with the density vertex at $\vect{q}\to0$ reproduces $\varepsilon_2^{\mathrm{PA}}$ and the measured edge, which tests everything except the kinetic vertex. Exciton final states near the edge enter as for the density row. With one momentum in place of the 630 targets of the density row, the row costs a small fraction of that row, and it enters \cref{sec:10.7-scalar,sec:13.8} and the scalar row of \cref{sec:21.7}.

\subsection{The \texorpdfstring{$f$}{f}-sum check with its parts separated}

The sum rule \cref{eq:E.7-fsum} tests whether the intermediate-state sum is complete. With norm-conserving pseudopotentials the band sum $\sum_m(\varepsilon_{m,\vect{k}+\vect{q}} - \varepsilon_{v\vect{k}})|\rho_{mv}(\vect{k},\vect{q})|^2$ at $\vect{k} = \mathrm{\Gamma}$ reaches 0.81 of $q^2/2m_e$ with 16 bands and 0.89 with 100 bands, at $|\vect{q}| = 0.1$ to \SI{0.5}{\angstrom^{-1}}, still rising slowly (\cref{sec:13.9}). That shortfall mixes the bands left out with the nonlocal term $\delta_{\mathrm{NL}}$, which a comparison with $q^2/2m_e$ alone ignores, and the two can be separated exactly.

The truncation part needs no more bands. Let $P^\perp_N$ project out the lowest $N$ eigenstates of $H(\vect{k}+\vect{q})$ and let $|\phi\rangle = P^\perp_N|u_{v\vect{k}}\rangle$. Since $\rho_{mv}(\vect{k},\vect{q}) = \langle u_{m,\vect{k}+\vect{q}}|u_{v\vect{k}}\rangle$ and $P^\perp_N$ commutes with $H(\vect{k}+\vect{q})$,
\begin{equation}
  \sum_{m>N}(\varepsilon_{m,\vect{k}+\vect{q}} - \varepsilon_{v\vect{k}})\,|\rho_{mv}(\vect{k},\vect{q})|^2 = \langle\phi|H(\vect{k}+\vect{q}) - \varepsilon_{v\vect{k}}|\phi\rangle ,
  \label{eq:E.14-remainder}
\end{equation}
one application of the Hamiltonian to a projected vector, the cost of a single eigensolver iteration. The nonlocal part is a ground-state expectation value, $\delta_{\mathrm{NL}} = \langle u_{v\vect{k}}|V_{\mathrm{NL}}(\vect{k}+\vect{q}) - V_{\mathrm{NL}}(\vect{k})|u_{v\vect{k}}\rangle$ minus its part linear in $\vect{q}$, from the same projectors as the Hamiltonian. The truncated sum plus \cref{eq:E.14-remainder} then equals $\vect{q}\cdot\vect{v}^{\mathrm{g}}_v + q^2/2m_e + \delta_{\mathrm{NL}}$ to the precision of the eigenvectors, a test of the code rather than of the physics. The ratio $\delta_{\mathrm{NL}}/(q^2/2m_e)$ against $q$ says how much of the 0.89 shortfall is the nonlocal term rather than missing bands; an all-electron (PAW) calculation, which has no nonlocal term (\cref{app:E.2}), checks it. The remainder also measures the intermediate-state weight above the Wannier manifold that the Sternheimer remainder supplies: with the nonlocal term, 60 bands should exhaust at least \SI{95}{\percent} of the sum. The check runs at $\mathrm{\Gamma}$ and five off-grid points, for the three cubic directions, in minutes.

The loss function and the rate replace the factor-ten estimate in the table of the window of \cref{sec:13.9}, and the reach replaces its curve.
